\documentclass{resume_en}
\usepackage[UTF8]{ctex} % Use ctex for Chinese support
\usepackage{cite}
\usepackage{hyperref}
\usepackage{bookmark}
\hypersetup{
colorlinks=true,
linkcolor=cyan,
filecolor=magenta,
urlcolor=blue,
}

\begin{document}
\pagenumbering{gobble}

%***"%"后面的所有内容是注释而非代码，不会输出到最后的PDF中
%***使用本模板，只需要参照输出的PDF，在本文档的相应位置做简单替换即可
%***修改之后，输出更新后的PDF，只需要点击Overleaf中的“Recompile”按钮即可

%在大括号内填写其他信息，最多填写4个，但是如果选择不填信息，
%那么大括号必须空着不写，而不能删除大括号。
%\otherInfo后面的四个大括号里的所有信息都会在一行输出
%如果想要写两行，那就用两次这个指令（\otherInfo{}{}{}{}）即可

%***********Personal Information**************
{\Huge\textbf{Fenghao Li}}

\vspace{0.5em}
\normalsize
Email: trisolaris03@gmail.com \quad | \quad Phone: +86 153-4335-6293 \quad | \quad GitHub: \href{https://github.com/s1amese2003}{github.com/s1amese2003}

%***********Education**************
\section{Education}

\datedsubsection{\textbf{Central South University of Forestry and Technology}, Information and Computing Science, \textit{Bachelor}}{2022.9 - 2026.6}

%***********Awards & Honors**************
\section{Competition Experience}

{\footnotesize
\datedsubsection{Bronze Medal, ICPC Asia Regional Contest}{}

\datedsubsection{Silver Medal, HNCPC Hunan Collegiate Programming Contest}{}
}

%***********Technical Skills**************
\section{Technical Skills}

{\footnotesize
\datedsubsection{Proficient in C++, familiar with STL standard library, concurrent programming, memory management, and performance optimization.}{}

\datedsubsection{Familiar with Redis internals, including data structure implementations, Sentinel model, persistence mechanisms, expiration policies, and memory eviction strategies.}{}

\datedsubsection{Familiar with MySQL fundamentals, including indexing, transaction isolation levels, locking mechanisms, and logging.}{}

\datedsubsection{Familiar with various network protocols such as TCP, IP, DHCP, OSPF, VLAN, STP, etc.}{}

\datedsubsection{Proficient in using Wireshark for packet capture and analyzing common network issues.}{}
}

%***********Project Experience**************
\section{Project Experience}

\datedsubsection{\textbf{C++ WeChat-like Instant Messaging System}}{}

\datedsubsection{\textbf{Instant Messaging \& Protocol Design}}{}
\Content
{Designed bidirectional persistent connection protocol based on TCP/WebSocket, using Protobuf for message structure definition, supporting text messages, file chunking, friend requests, and moments updates.}
{Implemented message routing and offline message queue with sequence number + ACK retransmission mechanism, improving reliability and consistency in weak network scenarios.}
{Tech Stack: C++, Boost.Asio, WebSocket, Protobuf, SQLite.}

\datedsubsection{\textbf{File Transfer \& Resumable Upload}}{}
\Content
{Designed chunked upload protocol with offset + chunk + MD5 validation, implementing resumable transfer for large files and concurrent chunk writing.}
{Server-side supports file persistence, task recovery, and streaming download with concurrent upload and integrity verification.}
{Tech Stack: C++, libcurl / cpp-httplib, Concurrent IO, MD5 Checksum.}

\datedsubsection{\textbf{Friend System \& State Management}}{}
\Content
{Implemented friend request/accept/delete/blacklist features, designed friend relationship state machine and authorization logic, with client-side SQLite caching of relationship snapshots.}
{Supports incremental synchronization and local queries, reducing network overhead and improving client responsiveness.}
{Tech Stack: C++, SQLite, ACL Policy.}

\datedsubsection{\textbf{Moments \& Multimedia}}{}
\Content
{Supports posting text and multiple images, server-side object storage for image metadata management, client-side thumbnail generation and caching strategy.}
{Designed fetch/incremental update interface with bandwidth control, optimizing loading and display experience on mobile networks.}
{Tech Stack: C++, HTTP File Service, OpenCV / Qt Image Processing, Caching Strategy.}

\datedsubsection{\textbf{Technical Highlights \& Engineering}}{}
\Content
{End-to-end implementation of instant messaging stack: protocol design, persistent connections, reliable transmission, file chunking, media signaling, and client UI.}
{Improved concurrent throughput through multi-threading and asynchronous IO (std::thread, mutex, Boost.Asio); containerized project (Docker) with reproducible demo.}
{Summary: C++/Qt + Boost.Asio WeChat-like IM implementation featuring text, file transfer, friend system, moments, and audio/video signaling.}


\end{document}